% Options for packages loaded elsewhere
\PassOptionsToPackage{unicode}{hyperref}
\PassOptionsToPackage{hyphens}{url}
\PassOptionsToPackage{dvipsnames,svgnames,x11names}{xcolor}
\documentclass[
  11pt,
  letterpaper,
  titlepage]{article}
\usepackage{xcolor}
\usepackage[margin=0.85in]{geometry}
\usepackage{amsmath,amssymb}
\setcounter{secnumdepth}{5}
\usepackage{iftex}
\ifPDFTeX
  \usepackage[T1]{fontenc}
  \usepackage[utf8]{inputenc}
  \usepackage{textcomp} % provide euro and other symbols
\else % if luatex or xetex
  \usepackage{unicode-math} % this also loads fontspec
  \defaultfontfeatures{Scale=MatchLowercase}
  \defaultfontfeatures[\rmfamily]{Ligatures=TeX,Scale=1}
\fi
\usepackage{lmodern}
\ifPDFTeX\else
  % xetex/luatex font selection
\fi
% Use upquote if available, for straight quotes in verbatim environments
\IfFileExists{upquote.sty}{\usepackage{upquote}}{}
\IfFileExists{microtype.sty}{% use microtype if available
  \usepackage[]{microtype}
  \UseMicrotypeSet[protrusion]{basicmath} % disable protrusion for tt fonts
}{}
\usepackage{setspace}
\makeatletter
\@ifundefined{KOMAClassName}{% if non-KOMA class
  \IfFileExists{parskip.sty}{%
    \usepackage{parskip}
  }{% else
    \setlength{\parindent}{0pt}
    \setlength{\parskip}{6pt plus 2pt minus 1pt}}
}{% if KOMA class
  \KOMAoptions{parskip=half}}
\makeatother
\setlength{\emergencystretch}{3em} % prevent overfull lines
\providecommand{\tightlist}{%
  \setlength{\itemsep}{0pt}\setlength{\parskip}{0pt}}
\usepackage{bookmark}
\IfFileExists{xurl.sty}{\usepackage{xurl}}{} % add URL line breaks if available
\urlstyle{same}
% fallback for those not using the hyperref driver hyperxmp:
\makeatletter
\@ifundefined{xmpquote}{\newcommand{\xmpquote}[1]{#1}}{}
\makeatother
\hypersetup{
  pdftitle={TentBag-Reworked Player Guide},
  colorlinks=true,
  linkcolor={blue},
  filecolor={Maroon},
  citecolor={Blue},
  urlcolor={blue},
  pdfcreator={LaTeX via pandoc}}

\title{TentBag-Reworked Player Guide}
\usepackage{etoolbox}
\makeatletter
\providecommand{\subtitle}[1]{% add subtitle to \maketitle
  \apptocmd{\@title}{\par {\large #1 \par}}{}{}
}
\makeatother
\subtitle{TentBag-Reworked 3.2.4 for Vintage Story 1.22.6}
\author{}
\date{2026-08-14}

\begin{document}
\maketitle

{
\hypersetup{linkcolor=black}
\setcounter{tocdepth}{2}
\tableofcontents
}
\setstretch{1.08}
TentBag-Reworked adds a craftable tent bag that can pack a small built
structure into an item, then place it again later. It is meant for
portable camps, tents, and other compact travel builds.

This guide describes the default behavior in TentBag-Reworked
\texttt{3.2.4} as indexed for Vintage Story \texttt{1.22.6}. Server
configuration can change the capture size, hunger cost, banned blocks,
waypoint behavior, and some placement rules.

\section{Crafting}\label{crafting}

Craft an Empty Tent Bag in the crafting grid with:

\begin{itemize}
\tightlist
\item
  \texttt{1} cloth in the center
\item
  \texttt{4} rope in the corners
\item
  \texttt{4} sticks on the side-center slots
\end{itemize}

The recipe also works with the rope and stick positions swapped. The
cloth color determines the bag color. Supported colors are plain, black,
blue, brown, gray, green, mordant, orange, pink, purple, red, white, and
yellow.

The tent bag does not stack.

\section{Bundled Sample Tent}\label{bundled-sample-tent}

TentBag-Reworked also ships a sample tent schematic named
\texttt{tent.json}. On server startup, the mod copies it to the game
data folder as \texttt{WorldEdit/tent.json} if that file does not
already exist.

This is a WorldEdit schematic/template, not a pre-packed Tent Bag and
not an extra survival recipe. Regular players do not automatically
receive it. In survival play, craft an Empty Tent Bag and pack a tent or
camp you built yourself.

If you have admin, creative, or WorldEdit access, you can use the copied
\texttt{tent.json} as a starter/example tent, then pack that structure
into a tent bag like any other build. The bundled schematic is
\texttt{7\ x\ 3\ x\ 6} and contains a small camp setup with blocks such
as a hay bed, firepit, storage vessels, barrel, chest, trunk, tool rack,
support beams, and ground storage.

\section{Using The Starter Tent}\label{using-the-starter-tent}

You must be in creative mode and have the \texttt{worldedit} privilege.
The file is loaded by filename from the game or server data
\texttt{WorldEdit} folder, so use \texttt{tent.json}, not an asset path.

For a direct server-side import:

\begin{enumerate}
\def\labelenumi{\arabic{enumi}.}
\tightlist
\item
  Stand at the point where you want the bottom center of the tent.
\item
  Run \texttt{/we\ start}, then run \texttt{/we\ end} without moving.
  This gives WorldEdit both ends of a tiny selection and creates the
  start marker that import needs.
\item
  Optional: run \texttt{/we\ import-rotation\ 90} or
  \texttt{/we\ import-rotation\ 180} before importing.
\item
  Run \texttt{/we\ import\ tent.json\ BottomCenter}.
\item
  If it lands wrong, run \texttt{/we\ undo}.
\end{enumerate}

\texttt{BottomCenter} centers the \texttt{7\ x\ 3\ x\ 6} schematic
horizontally on your marker. Without it, WorldEdit uses the marker as
the schematic's start corner, which is harder to place by eye.

Use the capitalization \texttt{BottomCenter}; lower-case input falls
back to the default corner placement.

For a preview workflow:

\begin{enumerate}
\def\labelenumi{\arabic{enumi}.}
\tightlist
\item
  Select the import tool with \texttt{/we\ tool\ import} or
  \texttt{/we\ t\ import}.
\item
  Load the schematic with \texttt{/we\ ims\ tent.json}.
\item
  Use \texttt{/we\ imr\ 90} to rotate the preview, or
  \texttt{/we\ imflip} to flip it, if needed.
\item
  Move the preview with \texttt{/we\ north}, \texttt{/we\ south},
  \texttt{/we\ east}, \texttt{/we\ west}, \texttt{/we\ up}, or
  \texttt{/we\ down}.
\item
  Right-click a block to place the preview, or run \texttt{/we\ apply},
  \texttt{/we\ place}, or \texttt{/we\ commit}.
\end{enumerate}

After the starter tent is pasted into the world, it is just normal
blocks. To turn it into a Tent Bag, craft or spawn an Empty Tent Bag,
aim at the top face of the ground block under the center of the tent,
and right-click to pack it.

\section{Packing A Tent}\label{packing-a-tent}

\begin{enumerate}
\def\labelenumi{\arabic{enumi}.}
\tightlist
\item
  Build the camp or tent you want to carry.
\item
  Hold an Empty Tent Bag.
\item
  Aim at the top face of the ground block under the center of the build.
\item
  Right-click to pack the area.
\end{enumerate}

By default, the packed area is a \texttt{7\ x\ 7\ x\ 7} box centered on
the clicked block, starting above the clicked ground block. The floor
block itself is not packed unless the server enables floor grabbing.

If packing succeeds, the empty bag becomes a Packed Tent Bag and the
original blocks are removed. In default configuration, the packed bag
stays in your active slot. If the server enables copy mode, the original
blocks are copied instead of removed.

Packing consumes satiety in survival mode. The default cost is
\texttt{2} satiety per packed non-air block. If you do not have enough
saturation, the mod refuses to pack and shows:
\texttt{I\textquotesingle{}m\ too\ hungry\ to\ pack\ this\ right\ now!}

\section{Placing The Tent}\label{placing-the-tent}

\begin{enumerate}
\def\labelenumi{\arabic{enumi}.}
\tightlist
\item
  Hold a Packed Tent Bag.
\item
  Aim at the ground where you want the tent centered.
\item
  Use the preview and rotate keys if needed.
\item
  Right-click to unpack.
\end{enumerate}

The packed build is placed above the block you click. The Packed Tent
Bag changes back into an Empty Tent Bag after unpacking.

Unpacking also consumes satiety in survival mode, using the same saved
block count from when the tent was packed. If you do not have enough
saturation, the mod refuses to unpack and shows:
\texttt{I\textquotesingle{}m\ too\ hungry\ to\ unpack\ this\ right\ now!}

By default, unpacking creates a map waypoint for the tent, and packing
removes that waypoint.

\section{Hotkeys}\label{hotkeys}

The item tooltip shows the current keybindings. Defaults are:

\begin{itemize}
\tightlist
\item
  \texttt{B}: toggle the packed tent hologram preview
\item
  \texttt{P}: toggle the empty bag capture-area preview
\item
  \texttt{R}: rotate a packed tent by \texttt{90} degrees
\end{itemize}

The packed tent preview follows the block you are aiming at. If
placement would fail, the mod highlights problem blocks.

\section{What Cannot Be Packed}\label{what-cannot-be-packed}

The default banned list prevents packing some world, resource, food, and
utility blocks. Common examples include:

\begin{itemize}
\tightlist
\item
  water and lava
\item
  crops, herbs, mushrooms, and berry bushes
\item
  farmland, raw clay, peat, rocks, and ores
\item
  grown logs and resin logs
\item
  static translocators and teleporter bases
\item
  paper lanterns and chandeliers
\item
  anvils and forges
\item
  cooked crocks, bowls of meal, and cooked clay pots
\end{itemize}

Servers can change this list or switch to an allow-list mode.

\section{Troubleshooting}\label{troubleshooting}

\texttt{The\ tent\ can\ only\ be\ packed\ by\ clicking\ on\ the\ ground!}

You tried to pack by clicking the side or bottom of a block. Aim at the
top face of the ground block under the build.

\texttt{There\ are\ no\ solid\ blocks\ in\ the\ build\ area!}

The capture box did not contain anything packable. Move the bag target
closer to the build center or check the server's configured capture
size.

\texttt{Cannot\ pack\ ...\ into\ the\ bag!}

The capture area contains a banned block. Remove that block or ask the
server owner whether the banned list is intentional.

\texttt{Cannot\ unpack\ here,\ a\ clear\ area\ is\ required!}

Something blocks the placement area. Clear the space or move the target.

\texttt{Cannot\ unpack\ here,\ solid\ ground\ is\ required!}

The server requires solid ground under the tent, and the target area
does not meet that rule.

\texttt{Failed\ to\ unpack:\ Internal\ error\ while\ reading\ saved\ data!}

The packed tent data could not be read. This usually means the item data
is damaged or incompatible.

\section{Notes For Players On
Servers}\label{notes-for-players-on-servers}

If TentBag-Reworked behaves differently from this guide, check with the
server owner. The server can change:

\begin{itemize}
\tightlist
\item
  capture radius and height
\item
  whether the floor is packed
\item
  whether packing removes blocks or only copies them
\item
  whether plants and loose rocks are replaced during placement
\item
  whether the bag stays in your inventory or drops as an item after use
\item
  waypoint creation
\item
  blocked or allowed block lists
\item
  satiety cost per packed block
\end{itemize}

\section{Reference Sources}\label{reference-sources}

This guide was written from the local indexed TentBag-Reworked
\texttt{3.2.4} reference files.

Under \texttt{versions/1.22.6/mods/}, use the
\texttt{tentbagreworked\_3.2.4/} mod root. Its asset files are in
\texttt{assets/tentbagreworked/}:

\begin{itemize}
\tightlist
\item
  \texttt{itemtypes/tentbagreworked.json}
\item
  \texttt{recipes/grid/tentbagreworked.json}
\item
  \texttt{config/tent.json}
\item
  \texttt{lang/en.json}
\end{itemize}

Under the same mod root, decompiled sources in
\texttt{decompiled\ assembly/TentBagReworked/}, then
\texttt{TentBagReworked/}:

\begin{itemize}
\tightlist
\item
  \texttt{TentBagReworkedModSystem.cs}
\item
  \texttt{Behaviors/PackableBehavior.cs}
\item
  \texttt{TentHologramRenderer.cs}
\item
  \texttt{Config/TentBagReworkedConfig.cs}
\item
  \texttt{Util/TentArea.cs}
\end{itemize}

Under the API root \texttt{versions/1.22.6/API/}, use the
\texttt{VSCreativeMod} decompile for version \texttt{1.22.6}, then
\texttt{Vintagestory/ServerMods/WorldEdit/}:

\begin{itemize}
\tightlist
\item
  \texttt{WorldEdit.cs}
\item
  \texttt{WorldEditWorkspace.cs}
\item
  \texttt{ImportTool.cs}
\end{itemize}

Under the same API root, use the \texttt{VintagestoryAPI} decompile for
version \texttt{1.22.6}, then \texttt{Vintagestory/API/Common/}:

\begin{itemize}
\tightlist
\item
  \texttt{BlockSchematic.cs}
\item
  \texttt{EnumOrigin.cs}
\end{itemize}

\end{document}
